% TOP_PROBLEM: {"id": 2832, "title": "Kirby Problem 3.34", "category_id": 7, "category": {"id": 7, "name": "topology", "display_name": "Topology", "description": "Properties preserved under continuous deformations.", "slug": "topology", "order_index": 7, "created_at": "2026-07-31T15:26:25.671Z"}, "status": "open", "problem_number": "KP-3.34", "set_id": 11}
% FIRST_POSTED: 2026-10-01
% ENTRY_KIND: statement_only
% UnsolvedMath ID 2832; source catalogue code KP-3.34
% !TeX program = lualatex
% Definition and sources only; this file is not an LLM solution attempt.
\documentclass[11pt,a4paper]{article}
\usepackage[margin=25mm]{geometry}
\usepackage{fontspec}
\setmainfont{lmroman10-regular.otf}[BoldFont=lmroman10-bold.otf,
 ItalicFont=lmroman10-italic.otf,BoldItalicFont=lmroman10-bolditalic.otf]
\usepackage{amsmath,amssymb,mathtools}
\usepackage{unicode-math}
\setmathfont{latinmodern-math.otf}
\AtBeginDocument{\let\setminus\smallsetminus}
\usepackage{xurl,hyperref}
\hypersetup{colorlinks=true,urlcolor=blue}
\setcounter{secnumdepth}{0}
\setlength{\parindent}{0pt}
\setlength{\parskip}{5pt}
\setlength{\emergencystretch}{3em}
\begin{document}
\section*{Kirby Problem 3.34}\label{2832}
{\small UnsolvedMath ID 2832 · KP-3.34 · Topology · Kirby's Problems in Low-Dimensional Topology}
\subsection{Definitions and mathematical statement}
A Heegaard splitting of a closed connected orientable three-manifold $M$
is a decomposition into two handlebodies meeting along their common
boundary surface. Its genus is the genus of that surface. A stabilization
is the connected sum of the splitting with the standard genus-one splitting
of $S^3$, performed in a ball. It raises genus by one without changing $M$.

Take any two genus-$g$ Heegaard splittings of $M$. Are they equivalent
after at most $g$ stabilizations of each? Equivalently, does there exist a
common stabilization of genus at most $2g$? Equivalence means ambient
isotopy of the splitting surface, allowing the two handlebodies to be
interchanged; a version remembering their order should be stated separately.

The quantifier includes splittings that are not of minimum genus for $M$.
One may choose the stabilization arcs and isotopies freely, but must obtain
a common stabilized splitting of the original manifold. The general
stabilization theorem guarantees a finite common stabilization, without
giving the particular bound $2g$ asked for here. Conversely, examples
requiring many stabilizations must be compared using the same equivalence
convention, since remembering which handlebody is first can affect the
required number.
\subsection{Short English statement}
Do two Heegaard splittings of the same genus always become isotopic after adding at most that many stabilizing handles to each?
\subsection{Sources}
\begin{itemize}
\item[S1] ULAM.AI and the UnsolvedMath Contributors, supplied problems.json, record 2832 (KP-3.34), Kirby's Problems in Low-Dimensional Topology. Catalogue identity and source transcription; CC BY 4.0.
\url{https://huggingface.co/datasets/ulamai/UnsolvedMath}.
\item[S2] K3: A New Problem List in Low-Dimensional Topology, Problem 3.34. Expanded formulation of the supplied catalogue statement.
\url{https://math.berkeley.edu/sites/default/files/surv-295-ruberman-watermarked-author-pdf.pdf}.
\end{itemize}
\end{document}
